\scanpage{13}
\textbf{C. R. Acad. Sc. Paris, t. 279 (16 septembre 1974)}\hfill\textbf{Série A — 449}

\rule{\linewidth}{0.4pt}

Pour traiter le cas général, rappelons que vu l'hypothèse faite sur A, les Algèbres $\mathrm A\otimes\mathrm A^*$
et $\operatorname{END}(\mathrm A)$ sont isomorphes. Par composition nous obtenons deux isomorphismes

\diagram{\begin{tikzcd}[ampersand replacement=\&]
\mathrm B\otimes\mathrm A\arrow[r,shift left=0.5ex,"{f\otimes1}"]\arrow[r,shift right=0.5ex,"{g\otimes1}"']\&
\mathrm A\otimes\mathrm A^*\arrow[r,"\tau"]\&\operatorname{END}(\mathrm A).
\end{tikzcd}}

Il existe donc $\bar s\in\mathrm U[(\mathrm A\otimes\mathrm A^*)_0]$ tel que
\[
\bar s^{-1}(g(b)\otimes a^*)\bar s=f(b)\otimes a^*.
\]
En donnant à $b$ la valeur 1, on voit que
\[
1\otimes a^*.\bar s=\bar s.1\otimes a^*.
\]
Mais nous savons que
\[
\mathrm C_{\mathrm A^*}(1\otimes\mathrm A^*)=1\otimes\mathrm Z(\mathrm A^*)=\mathrm A\otimes k=\mathrm A\otimes1
\]
ce qui montre que $\bar s=s\otimes1$ où $s\in\mathrm A_0$. Nous voyons que $\bar s^{-1}=t\otimes1$ et par conséquent
$s\in\mathrm U(\mathrm A_0)$. En décomposant $\bar s$, en faisant $a=1$, on a le résultat souhaité compte tenu de
\[
\mathrm A\simeq\mathrm A\otimes1\subseteq\mathrm A\otimes\mathrm A^*=\mathrm A^e.
\]

\textsc{Théorème 2}. — \emph{Soit} R \emph{un anneau commutatif semi-local,} A \emph{et} B \emph{deux R-Algèbres d'Azumaya.
Si $f$ et $g$ sont deux homomorphismes d'Algèbres injectifs de} B \emph{dans} A, \emph{il existe un
automorphisme intérieur $\theta$ induit par un élément de degré} 0 \emph{de} A \emph{tel que $f=\theta\circ g$.}

\emph{Preuve}. — Comme dans (${}^2$), on commence par démontrer ce résultat pour un produit
fini de corps, puis on considère les homomorphismes induits $\bar f$, $\bar g$ de $\mathrm B/\operatorname{Rad}(\mathrm R)\,\mathrm B$ dans
$\mathrm A/\operatorname{Rad}(\mathrm R)\,\mathrm A$. L'automorphisme intérieur $\bar\theta$ de $\mathrm A/\operatorname{Rad}(\mathrm R)\,\mathrm A$ se remonte en un automorphisme
$\theta$ de A.

3. \textsc{Régularité des algèbres d'Azumaya graduées}.

\textsc{Théorème 3}. — \emph{Soit} R \emph{un anneau semi-local,} A, B \emph{et} C \emph{des R-algèbres d'Azumaya telles
qu'il existe un isomorphisme $h$ entre $\mathrm A\otimes\mathrm C$ et $\mathrm B\otimes\mathrm C$. Alors il existe un isomorphisme
entre} A \emph{et} B.

\emph{Preuve}. — Il suffit de considérer les deux homomorphismes

\diagram{\begin{tikzcd}[ampersand replacement=\&]
\mathrm C\arrow[r]\arrow[d]\&1\otimes\mathrm C\arrow[r]\&\mathrm A\otimes\mathrm C\arrow[d]\\
1\otimes\mathrm C\arrow[rr]\&\&\mathrm B\otimes\mathrm C
\end{tikzcd}}

et d'appliquer le théorème de Skolem-Nother.

(*) Séance du 24 juin 1974.

(${}^1$) N. \textsc{Bourbaki}, \emph{Modules et anneaux semi-simples}, Hermann, Paris, 1958.

(${}^2$) A. \textsc{Roy} et R. \textsc{Sridharan}, \emph{J. Math. Kyoto Univ.}, \textbf{7-2}, 1967, p. 161-167.

(${}^3$) C. \textsc{Small}, \emph{Trans. Amer. Math. Soc.}, \textbf{156}, 1971, p. 455-491.

\begin{flushright}
\emph{Université des Sciences}\\
\emph{et des Techniques du Languedoc,}\\
\emph{U. E. R. de Mathématiques,}\\
\emph{place Eugène-Bataillon,}\\
34060 \emph{Montpellier.}
\end{flushright}

\mspage{12}
